\section{Asuka Langley: tsundere, pride, injury, performance, fear of worthlessness}

Аска — наш самый часто цитируемый персонаж (1.4, 3.3, 5.2, 10.3), и здесь стоит не пересказывать уже разобранный механизм превосходство/исчезновение, а зайти с ракурса, которого мы ещё не касались: страх никчёмности как источник, а не следствие её характера.

Вектор: affect — взрывной поверх тревожного; роль — соперница, стремящаяся к признанию через результат; intimacy — контролирующая, близкая к паттерну Сенджогахары (5.1, 11.6), но менее осознанная, более защитная; agency — изначально сильная, но хрупкая, потому что зависит от постоянного подтверждения результатом (5.7, ось self-worth); conflict — атака и превосходство как основной режим, исчезновение как резервный.

Деталь, добавляющая объём: её пилотирование EVA и её самоценность связаны настолько тесно, что технический сбой систем в кризисный момент истории читается ею не как техническая проблема, а как экзистенциальное доказательство её ненужности, — то есть её self-worth (5.7) буквально неотделим от результата работы, и это ровно тот failure mode, от которого 5.7 предостерегает применительно к Рин.

Для книги Аска — предупреждение о том, что бывает, когда self-worth полностью завязан на полезность: не карикатурная грубость, а постоянная, изматывающая необходимость доказывать право на существование через результат.
